\section{Appendix: Scope and Qualifications}
This concise study derives from \work{The Word, the Wound, and the Song},
the same canonical owner's expansive study. It treats the 1962 Missal's
Twentieth Sunday after Pentecost, second class, green, on 11 October
2026 under the general calendar. The FSSP Province of France's October
Ordo corroborates the Sunday within its institutional scope. No local
patronal or dedication calendar was supplied; current permission to
celebrate in a particular church is not determined here. The three
Maternity prayers are conditional as stated in the opening map.

The two whole-formulary readings are editorial syntheses over the
checked commentaries, not works attributed to a Father. Schuster's
matching Sunday exposition does not cover the dated Marian commemoration
or the Trinity Preface. Exact text, evidence and interpretation records
belong to \texttt{propers/verified.md}, \texttt{research/scope.md} and
\texttt{research/interpretations.md}. The reception search is bounded
by the cited loci. Chronological alternatives remain historical scholarly
judgments: an author's reference era does not date a psalm's occasion;
the Psalter boundary does not securely date each poem. The Petavius
table retains its printed-page verification limit.

This edition omits the complete appointed texts. Short biblical wording
uses the historical Douay--Rheims--Challoner witness; chant adaptations
are distinguished where they affect interpretation. Cummiskey and
Lasance--Walsh provide the expansive study's historical prayer English,
on their recorded United States public-domain bases. It is study material,
not an approved English liturgical edition. No generated prayer translation
or restricted modern liturgical text is supplied. This is an AI-assisted
Catholic study aid, not an official liturgical book or ecclesiastical judgment.

\section{References}
\begin{description}[style=nextline,leftmargin=0pt,labelsep=0pt,itemsep=.35em]
\item[Appointments and text]
\work{Missale Romanum}, Vatican, 1962, pp.~412--413, 685--686,
244--245; General Rubrics 91,106--109,111; Mass Rubrics
433--437,480,483,494,505. CMAA facsimile:
\url{https://media.churchmusicassociation.org/pdf/missale62.pdf}.
FSSP France, \work{Ordo du mois}, October 2026, 11 October row,
consulted 5 October 2026, \url{https://www.fssp.fr/ordo-du-mois/}.
Douay--Rheims--Challoner, Gutenberg eBook 1581: Daniel~3;
Psalms~107,118,136,144; John~4; Ephesians~4--6; Jeremiah~25.

\item[Patristic and saintly commentary]
Augustine, \work{Tractates on John} XVI.1--7 (NPNF I.7);
\work{Expositions of the Psalms}, Ps.~107.1--2; 118.50--56;
136.2--13; 144.14--17 (NPNF I.8, Hebrew-numbered headings).
Gregory, \work{Homilies on the Gospels}, II.28.1--3, Latin
Wikisource transcription, contributors including Mizardellorsa,
CC BY-SA 3.0; the present account paraphrases, not translates.
Chrysostom, \work{Homilies on Ephesians} XIX (NPNF I.13), with
editorial note~404 separately identified. Jerome, \work{Commentary
on Daniel} I, at 3:23--29, PL~25, cols.~639--640.
Bellarmine, \work{Commentary on the Psalms}, O'Sullivan translation
(1866), Ps.~136:1--4,144:14--17. Schuster, \work{The Sacramentary},
Levelis-Marke translation, III (1927), pp.~174--178.

\item[Historical attribution and chronology]
Haydock, \work{Douay--Rheims Bible with Commentary}, Loreto memorial
edition (2014), notes at Ps.~70 (p.~740), 118:1 (p.~781),136 (p.~794).
\work{Catholic Encyclopedia}: Corbett, ``David, King''; Gigot,
``Book of Daniel,'' deuterocanonical portions; Ladeuze, ``Epistle to
the Ephesians,'' date/place; Prat, ``St. Paul,'' chronology; Fonck,
``Gospel of St. John,'' date; Durand, ``New Testament,'' Johannine
date; Maas, ``Chronology of the Life of Jesus Christ,'' second journey;
Reid, ``Captivities of the Israelites''; Meistermann, ``Temple of
Jerusalem''; Schets, ``Third and Fourth Books of Kings''; Sloet,
``Chronology of the Kings,'' Petavius table. USCCB NABRE introductions,
Psalms and Ephesians, consulted 5 October 2026:
\url{https://bible.usccb.org/bible/psalms/0};
\url{https://bible.usccb.org/bible/ephesians/0}.

\item[Historical English prayers]
Cummiskey, \work{Roman Missal for the Use of the Laity} (1861),
pp.~448--450, xxviii--xxix. Lasance--Walsh, \work{The New Roman
Missal}, revised 1945 edition, later printing, pp.~1233,1235--1236.
Historical texts retain their own rights; the project license covers
the original exposition, not its sources.
\end{description}
